%!TEX program = lualatex

\documentclass[portuguese, 11pt, a4paper, titlepage, oneside]{article}

\usepackage{geometry}
\usepackage[T1]{fontenc}
\usepackage{graphicx}
\usepackage{parskip}
\usepackage{amssymb}
\usepackage[portuguese]{babel}
\usepackage[framed, numbered]{matlab-prettifier}
\usepackage{subcaption}
\usepackage{wrapfig}
\usepackage{hyperref}

% Title page config
\title{
  \includegraphics[width=10cm]{assets/deti_horizontal_preto.pdf}
  \vskip 2em
  \huge Aspetos Profissionais e Sociais\\
  da Engenharia Informática
  \vskip 1em
	{\bfseries Trabalho individual 1}
}

\author{Rúben Garrido\\NMec: 107927}
\date{\today}
% Begin document
\begin{document}

% Title
\maketitle

\section{Utilização do ChatGPT pelos alunos e professores, no contexto de APSEI}
\subsection{Introdução}
A inteligência artificial baseada em \textit{Large Language Models} (neste caso, o ChatGPT) revolucionou a forma como os humanos pensam e interagem com a tecnologia. Todas as áreas foram influenciadas, embora a educação tenha tido um dos maiores destaques.

Esta mudança de paradigma deve-se sobretudo à utilização do ChatGPT (bem como de outros LLMs) pelos alunos, por forma a "facilitar" a curva de pensamento necessária para realizar um determinado trabalho. No entanto, o uso excessivo e irresponsável da inteligência artificial levou ao surgimento de respostas meramente criadas por esta, com imprecisões e padrões repetidos.

Aspetos Profissionais e Sociais da Engenharia Informática é uma Unidade Curricular que, ao contrário das restantes UCs do curso, se foca em questões éticas e de posicionamento dos engenheiros informáticos perante a sociedade.

\subsection{Alunos}
A meu ver, a utilização do ChatGPT pelos alunos nesta unidade curricular \textbf{deve ser permitida, e até expectável} por parte dos professores. No entanto, e como tudo, deve ser utilizada com moderação e com espírito crítico.
\subsubsection{Aspetos positivos}
A integração da inteligência artificial vem trazer algumas ideias novas, em situações onde nem sempre se consegue discernir por nós próprios. Esta permite simular alguns dilemas éticos, e inclusive olhar para um certo problema com uma visão diferente. [\textit{Dizer que estou a utilizar o Gemini para pensar em tópicos}]

Para além disso, a utilização do ChatGPT pode ser benéfica \textbf{no que toca à correção linguística, gramatical e ortográfica}. Uma vez que APSEI é a única unidade curricular do curso que onde é necessário realizar trabalhos de escrita, é comum os alunos não se sentirem tão confortáveis a escrever artigos de opinião, devido à falta de prática.

Assim, e por um lado, o ChatGPT permite corrigir erros gramaticais e reescrever frases menos bem construídas. Por outro, quando utilizada com espírito crítico, pode ser um incentivo aos alunos de modo a que estes desenvolvam \textit{skills} linguísticas mais aprofundadas.

\subsubsection{Aspetos negativos}
Texto

\subsection{Professores}
Texto

\section{O que falta abordar}
\begin{itemize}
  \item If yes, should students be expected to use ChatGPT?
  \item If yes, should students be incentivized to use ChatGPT in which way?
  \item If no, why not?
  \item And can/should professors follow the same rules as the students?
\end{itemize}
\end{document}